% ====================================================================
\chapter{Mathematical Induction}\label{ch:ind}
% ====================================================================

\begin{chapterintro}
Many statements in mathematics and economics are claims about every natural number at once: a formula for the sum of
the first $n$ numbers, a bound on a sequence that holds in every period, a property of a polynomial of every degree.
Checking such a claim for a few values of $n$ proves nothing about the infinitely many values that remain.
Mathematical induction is the technique that proves all of them together, by proving just two things.
\end{chapterintro}

\begin{prerequisites}
Propositions and implication (\cref{ch:logic}); the natural numbers and the successor property (\cref{sec:number-systems});
summation notation and property P3 (\cref{ch:sum}).
\end{prerequisites}

\section{The principle of induction}\label{sec:ind-principle}

Induction applies when there is a sequence of similar propositions $P_1,P_2,\dots,P_n,\dots$ that depend on a
parameter $n\in\N$. \src{L1, slide 32; L2, slide 4}

\begin{principle}[Mathematical induction][pr:induction]
Instead of proving each of the propositions $P_1,P_2,P_3,\dots$ separately, we prove only two propositions:
\begin{itemize}
  \item the \term{induction base}: the first proposition, $P_1$;
  \item the \term{induction step}: the proposition that for any index $n$, $P_n\Rightarrow P_{n+1}$.
\end{itemize}
When both the induction base and the induction step are proven, all the propositions $P_1,P_2,P_3,\dots$ are proven.
When proving the induction step $P_n\Rightarrow P_{n+1}$, the statement that $P_n$ is true is called the
\term{induction hypothesis}. \src{L2, slide 4}
\end{principle}

\begin{intuition}[Dominoes]
Imagine an infinite row of dominoes, one for each natural number. The base knocks over the first domino. The step
guarantees that whenever a domino falls, it knocks over the next one. Then every domino falls:
\[
P_1\ \Longrightarrow\ P_2\ \Longrightarrow\ P_3\ \Longrightarrow\ P_4\ \Longrightarrow\ P_5\ \Longrightarrow\ \cdots
\]
Both parts are essential. A perfect chain of dominoes achieves nothing if nobody pushes the first; and pushing the
first achieves nothing if there is a gap somewhere in the row.
\end{intuition}

\begin{figure}[ht]
\centering
\begin{tikzpicture}[x=11mm]
  \foreach \i in {1,...,7} {
    \pgfmathsetmacro{\ang}{\i<4 ? -58 : 0}
    \begin{scope}[shift={(\i,0)}]
      \fill[accent!\ifnum\i<4 70\else 25\fi, rotate around={\ang:(0.18,0)}] (0,0) rectangle (0.18,1.1);
      \draw[rotate around={\ang:(0.18,0)}] (0,0) rectangle (0.18,1.1);
      \node[below, font=\small] at (0.09,-0.05) {$P_{\i}$};
    \end{scope}
  }
  \node[font=\small] at (8.4,0.4) {$\cdots$};
  \draw[thick] (0.5,0) -- (8.8,0);
  \draw[-{Stealth}, thick, corecol] (0.2,1.2) -- (0.95,0.95) node[pos=0, above left, font=\small] {base};
  \draw[-{Stealth}, thick, excol] (3.7,1.5) to[bend left=30] (4.3,1.3);
  \node[font=\small, excol, above] at (4.0,1.6) {step: $P_n\Rightarrow P_{n+1}$};
\end{tikzpicture}
\caption{Induction as a row of dominoes. Proving the base $P_1$ knocks over the first; proving the step shows that each
falling domino knocks over its successor. Together they guarantee that every $P_n$ is true.}
\label{fig:dominoes}
\end{figure}

\begin{external}[Why induction is valid]
Suppose the base and the step have been proved, but some $P_n$ is false. Among the natural numbers $n$ for which
$P_n$ is false there is a smallest one, say $m$. It is not $1$, because $P_1$ is true; so $m-1\in\N$, and $P_{m-1}$ is true
because $m$ was the smallest failure. But the step gives $P_{m-1}\Rightarrow P_m$, so $P_m$ is true, a contradiction.
(The argument uses the fact that every non-empty set of natural numbers has a smallest element, which is
equivalent to the principle of induction itself. In this module induction is taken as a basic property of $\N$,
resting on the successor property.)
\end{external}

\section{The first example}\label{sec:ind-first}

\begin{example}[Sum of the first $n$ natural numbers][ex:sum-n]
Prove that for every $n\in\N$, $\displaystyle\sum_{j=1}^{n}j=\frac{n(n+1)}{2}$. \src{L1, slides 33--34; L2, slides 5--6}
\solution
For $n=1,2,3,\dots$ let $P_n$ be the proposition
\[
P_n:\qquad \sum_{j=1}^{n}j=\frac{n(n+1)}{2}.
\]
Put differently, $P_n$ states that the sum of all natural numbers from $1$ to $n$ is equal to $\frac{n(n+1)}2$.

\emph{Induction base.} The first proposition, $P_1$, is $\sum_{j=1}^{1}j=\frac{1\cdot(1+1)}2$. The left-hand side is $1$ and
the right-hand side is $1$, so $P_1$ is true.

\emph{Induction step.} We must show that for any natural $n$, $P_n\Rightarrow P_{n+1}$. Here $P_{n+1}$ is the proposition
\[
\sum_{j=1}^{n+1}j=\frac{(n+1)\big((n+1)+1\big)}{2}.
\]
Assume $P_n$ (the induction hypothesis). Then
\begin{align*}
\sum_{j=1}^{n+1}j&=\sum_{j=1}^{n}j+(n+1)=\frac{n(n+1)}{2}+(n+1)\\
&=\Big(\frac n2+1\Big)(n+1)=\frac{(n+2)(n+1)}{2}=\frac{(n+1)\big((n+1)+1\big)}{2},
\end{align*}
which is $P_{n+1}$. Note that we used $P_n$, the induction hypothesis, in the second equality.

\emph{Conclusion.} By the principle of induction, $P_n$ is true for every $n\in\N$.
\end{example}

The first equality in the step, $\sum_{j=1}^{n+1}j=\sum_{j=1}^{n}j+(n+1)$, is property P3 with $m=n+1$: we
\emph{peel off the last term}. This is the standard first move in any induction proof about a sum, because it exposes
the expression $\sum_{j=1}^n j$ about which the hypothesis tells us something.

\begin{keyformula}[Sum of the first $n$ natural numbers][kf:sum-n]{\sum_{j=1}^{n}j=1+2+\dots+n=\frac{n(n+1)}{2}}
\fsymbols $n\in\N$ is the number of terms; $j$ is the summation index.
\fmeaning The total of the first $n$ natural numbers. Equivalently, their average is $\frac{n+1}2$, the average of the
first and last terms.
\forigin Proved by induction in \cref{ex:sum-n}; recalled quickly by pairing the first and last terms
(\cref{sec:sigma-compute}).
\fuse To evaluate sums of linear expressions $\sum(aj+b)$ together with P1, P2 and the constant-sum rule; in
economics, for totals of quantities growing by a fixed amount each period.
\fexample $\sum_{j=1}^{100}j=\frac{100\cdot101}2=5050$. And $\sum_{j=1}^{12}(1000+50j)=12000+50\cdot78=15900$.
\fmistakes Applying the formula when the sum does not start at $1$; use P3: $\sum_{j=k}^{n}j=\frac{n(n+1)}2-\frac{(k-1)k}2$.
Writing $\frac{n(n-1)}2$, which is $\sum_{j=1}^{n-1}j$.
\frearrange Doubling: $2\sum_{j=1}^n j=n^2+n$. Shifting: $\sum_{j=0}^{n-1}j=\frac{(n-1)n}2$.
\fchanges Adding one more term adds $n+1$; the sum grows roughly like $\frac{n^2}2$, i.e.\ quadratically in $n$.
\fexam As the standard first example of induction, and as a given fact (``you may take as given that\dots'') in
computations with $\sum$.
\end{keyformula}

\section{The structure of an induction proof}\label{sec:ind-structure}

\begin{core}[Template for a proof by induction]
\begin{enumerate}[label=\arabic*.]
  \item \textbf{State $P_n$.} Write down precisely the proposition $P_n$ to be proved for each $n$.
  \item \textbf{Base.} Verify $P_1$, evaluating both sides separately.
  \item \textbf{Hypothesis.} Fix an arbitrary $n\in\N$ and assume $P_n$.
  \item \textbf{Target.} Write out $P_{n+1}$, so you know what you are aiming for.
  \item \textbf{Step.} Starting from one side of $P_{n+1}$, use the hypothesis (and say where) to reach the other side.
  \item \textbf{Conclusion.} ``By the principle of induction, $P_n$ holds for every $n\in\N$.''
\end{enumerate}
\end{core}

\begin{warning}[The hypothesis must be used]
If your argument for the step never uses $P_n$, then either the statement could have been proved directly (so
induction was unnecessary), or, far more often, there is an error. The hypothesis is the only extra information the
step is entitled to use.
\end{warning}

\begin{warning}[The base is not a formality]
The step alone proves nothing. Consider the false statement $P_n$: ``$n=n+1$''. Its step is perfectly valid: if
$n=n+1$, then adding $1$ gives $n+1=n+2$, which is $P_{n+1}$. But $P_1$ (``$1=2$'') is false, so no domino ever falls.
\end{warning}

\section{An equivalent form of the step}\label{sec:ind-alt-step}

The slides note that the step can equally be written with the indices shifted down by one.

\begin{example}[The step written as $P_{n-1}\Rightarrow P_n$]
Equivalently, the induction step is the proposition that for any natural $n\ge2$, $P_{n-1}\Rightarrow P_n$. In this
formulation the induction hypothesis is $P_{n-1}$. Carry out the step for \cref{ex:sum-n} in this form.
\src{L1, slide 35; L2, slide 7}
\solution
The hypothesis $P_{n-1}$ is $\sum_{j=1}^{n-1}j=\frac{(n-1)((n-1)+1)}2=\frac{(n-1)n}2$. Then
\[
\sum_{j=1}^{n}j=\sum_{j=1}^{n-1}j+n=\frac{(n-1)n}{2}+n=\Big(\frac{n-1}{2}+1\Big)n=\frac{(n+1)n}{2},
\]
so $P_n$ holds. Since the argument works for all $n\ge2$, $P_{n-1}\Rightarrow P_n$ for all $n\ge 2$, proving the step.
\end{example}

Both forms say the same thing: every proposition in the chain implies the next one. Use whichever makes the algebra
tidier.

\section{More examples}\label{sec:ind-more}

\begin{example}[Sum of the first $n$ odd numbers]
Prove that $\sum_{j=1}^{n}(2j-1)=n^2$ for every $n\in\N$.
\solution
Let $P_n$ be $\sum_{j=1}^n(2j-1)=n^2$.
\emph{Base:} for $n=1$ the left side is $2\cdot1-1=1$ and the right side is $1^2=1$.
\emph{Step:} assume $P_n$. Then
\[
\sum_{j=1}^{n+1}(2j-1)=\sum_{j=1}^{n}(2j-1)+\big(2(n+1)-1\big)=n^2+2n+1=(n+1)^2,
\]
using $P_n$ in the second equality. This is $P_{n+1}$, so by induction $P_n$ holds for all $n\in\N$.
\end{example}

\begin{example}[Sum of squares][ex:sum-squares]
Prove that $\displaystyle\sum_{j=1}^{n}j^2=\frac{n(n+1)(2n+1)}{6}$ for every $n\in\N$. \src{SG, mock Q5(a)}
\solution
Let $P_n$ be the displayed identity.
\emph{Base:} $n=1$ gives $1$ on the left and $\frac{1\cdot2\cdot3}{6}=1$ on the right.
\emph{Hypothesis:} assume $P_n$ for some $n\in\N$.
\emph{Target:} $\sum_{j=1}^{n+1}j^2=\frac{(n+1)(n+2)(2n+3)}6$.
\emph{Step:}
\begin{align*}
\sum_{j=1}^{n+1}j^2&=\sum_{j=1}^{n}j^2+(n+1)^2=\frac{n(n+1)(2n+1)}{6}+(n+1)^2 \qquad\text{(hypothesis used here)}\\
&=\frac{(n+1)\big[n(2n+1)+6(n+1)\big]}{6}=\frac{(n+1)(2n^2+7n+6)}{6}\\
&=\frac{(n+1)(n+2)(2n+3)}{6},
\end{align*}
since $2n^2+7n+6=(n+2)(2n+3)$. This is $P_{n+1}$, because $(n+1)+1=n+2$ and $2(n+1)+1=2n+3$. By induction the
identity holds for all $n\in\N$.
\end{example}

\begin{example}[A geometric sum][ex:geometric]
Let $r\ne1$ be a real number. Prove that for every $n\in\N$,
\[
\sum_{t=0}^{n-1}r^t=1+r+r^2+\dots+r^{n-1}=\frac{1-r^n}{1-r}.
\]
\solution
\emph{Base:} for $n=1$ the left side is $r^0=1$ and the right side is $\frac{1-r}{1-r}=1$.
\emph{Step:} assume the formula for $n$. Then
\[
\sum_{t=0}^{n}r^t=\frac{1-r^n}{1-r}+r^n=\frac{1-r^n+r^n-r^{n+1}}{1-r}=\frac{1-r^{n+1}}{1-r},
\]
which is the formula for $n+1$. By induction it holds for every $n\in\N$.
\end{example}

\begin{external}[Why the geometric sum matters in economics]
Present values of constant payment streams, annuity formulas and the Keynesian multiplier are all geometric sums.
For example, an asset paying $\pounds1$ at the end of each of $n$ years, at interest rate $i>0$, has present value
$\sum_{t=1}^{n}(1+i)^{-t}=\frac{1-(1+i)^{-n}}{i}$, obtained from \cref{ex:geometric} with $r=\frac{1}{1+i}$.
\end{external}

\subsection{Inequalities by induction, and other starting points}

Induction proves inequalities as well as identities, and the chain of dominoes need not start at $1$: to prove $P_n$
for all $n\ge n_0$, prove $P_{n_0}$ and $P_n\Rightarrow P_{n+1}$ for all $n\ge n_0$.

\begin{example}[Bernoulli's inequality][ex:bernoulli]
Let $x\ge-1$. Prove that $(1+x)^n\ge1+nx$ for every $n\in\N$. Deduce that $2^n\ge n+1$.
\solution
\emph{Base:} $(1+x)^1=1+x$, so equality holds.
\emph{Step:} assume $(1+x)^n\ge1+nx$. Since $x\ge-1$, the number $1+x$ is non-negative, so multiplying the hypothesis
by it keeps the direction (M6, or trivially if $1+x=0$):
\[
(1+x)^{n+1}=(1+x)^n(1+x)\ge(1+nx)(1+x)=1+(n+1)x+nx^2\ge1+(n+1)x,
\]
since $nx^2\ge0$. By induction the inequality holds for all $n\in\N$. With $x=1$: $2^n\ge1+n$.
\end{example}

\begin{example}[Starting from a later value]
Prove that $2^n>n^2$ for every integer $n\ge5$.
\solution
\emph{Base} ($n=5$): $32>25$.
\emph{Step:} assume $2^n>n^2$ for some $n\ge5$. Then $2^{n+1}=2\cdot2^n>2n^2$. It remains to show $2n^2\ge(n+1)^2$, i.e.\
$n^2-2n-1\ge0$, i.e.\ $(n-1)^2\ge2$, which holds for $n\ge3$. So $2^{n+1}>(n+1)^2$, and by induction the inequality holds
for all $n\ge5$. (It fails for $n=2,3,4$: $4=4$, $8<9$, $16=16$.)
\end{example}

\section{The alternative (strong) formulation}\label{sec:strong-ind}

\begin{principle}[Induction, alternative formulation][pr:strong]
To prove the propositions $P_1,P_2,\dots,P_n,\dots$:
\begin{itemize}
  \item the induction base is the first proposition $P_1$;
  \item the induction step is the proposition that for any index $n$, $(P_1\text{ and }\dots\text{ and }P_n)\Rightarrow P_{n+1}$.
\end{itemize}
When proving this step, the statement that the propositions $P_1,\dots,P_n$ are true is the induction hypothesis. The
two versions of the principle are, in fact, equivalent. \src{L1, slide 36; L2, slide 8}
\end{principle}

The hypothesis is now stronger --- all the earlier propositions, not only the last one --- so the step is easier to
prove. This is useful when $P_{n+1}$ depends on a case other than $P_n$, for example on two earlier cases, or on a case
with roughly half the size. The two formulations are equivalent: any proof using only $P_n$ is also a proof using all
of $P_1,\dots,P_n$; and conversely, applying ordinary induction to the statement $Q_n$: ``$P_1$ and $\dots$ and $P_n$''
recovers the strong form. The proof of the factorisation theorem for polynomials (\cref{thm:factor}) uses this
formulation.

\begin{example}[A recursively defined sequence][ex:strong-seq]
A sequence is defined by $x_1=1$, $x_2=3$ and $x_{n+1}=x_n+2x_{n-1}$ for $n\ge2$. Prove that $x_n=\dfrac{2^{n+1}+(-1)^n}{3}$
for every $n\in\N$.
\solution
Let $P_n$ be the formula for $x_n$. Since $x_{n+1}$ depends on \emph{two} earlier terms, we use the strong form and need
two base cases.
\emph{Base:} $n=1$: $\frac{4-1}{3}=1=x_1$. $n=2$: $\frac{8+1}{3}=3=x_2$. (The case $n=2$ must be checked separately
because the recursion only starts at $n\ge2$ and would require $x_0$, which is not defined.)
\emph{Step:} let $n\ge2$ and assume $P_1,\dots,P_n$; in particular $P_n$ and $P_{n-1}$. Then
\begin{align*}
x_{n+1}=x_n+2x_{n-1}&=\frac{2^{n+1}+(-1)^n}{3}+\frac{2\big(2^{n}+(-1)^{n-1}\big)}{3}\\
&=\frac{2\cdot 2^{n+1}+(-1)^n-2(-1)^{n}}{3}=\frac{2^{n+2}+(-1)^{n+1}}{3},
\end{align*}
using $(-1)^{n-1}=-(-1)^n$. This is $P_{n+1}$. By (strong) induction, $P_n$ holds for all $n\in\N$.
\end{example}

\section{What induction is not}\label{sec:ind-not}

\begin{warning}[Checking cases is not a proof]
A student verifies $\sum_{j=1}^{n}j^2=\frac{n(n+1)(2n+1)}6$ for $n=1,\dots,100$ and concludes that it holds for all $n$.
This establishes one hundred propositions and says nothing about $P_{101}$; by the successor property there is always a
further case. Only the general step, proved for an arbitrary $n$, covers them all.
\end{warning}

\begin{external}[A famous faulty induction]
``All horses have the same colour.'' Let $P_n$ be ``in any group of $n$ horses, all have the same colour''. $P_1$ is
true. For the step, take $n+1$ horses; remove one, and by the hypothesis the remaining $n$ have the same colour; remove a
different one, and again the remaining $n$ share a colour; the two groups overlap, so all $n+1$ share a colour.
The flaw is in the step from $n=1$ to $n=2$: the two groups of size $1$ do not overlap. The step must work for
\emph{every} $n$, including the smallest.
\end{external}

% --------------------------------------------------------------------
\section{Chapter review}
% --------------------------------------------------------------------
\subsection*{Summary}
Induction proves an infinite sequence of propositions $P_1,P_2,\dots$ by proving the base $P_1$ and the step
$P_n\Rightarrow P_{n+1}$ for arbitrary $n$; in the step, $P_n$ is the induction hypothesis and must be used. For sums, the
step starts by peeling off the last term (P3). The step may equally be written $P_{n-1}\Rightarrow P_n$ for $n\ge2$; the
chain may start at any $n_0$; and in the equivalent strong formulation the hypothesis is $P_1$ and $\dots$ and $P_n$, with
as many base cases as the argument needs. Checking finitely many cases is never a proof.

\subsection*{Key definitions}
\begin{itemize}
  \item \textbf{Induction base:} the proposition $P_1$.
  \item \textbf{Induction step:} for every $n$, $P_n\Rightarrow P_{n+1}$ (equivalently, $P_{n-1}\Rightarrow P_n$ for $n\ge2$).
  \item \textbf{Induction hypothesis:} the assumption that $P_n$ (or $P_1,\dots,P_n$) is true, made while proving the step.
\end{itemize}

\subsection*{Key equations}
\[
\sum_{j=1}^{n}j=\frac{n(n+1)}2,\qquad\sum_{j=1}^{n}(2j-1)=n^2,\qquad\sum_{j=1}^{n}j^2=\frac{n(n+1)(2n+1)}6,\qquad
\sum_{t=0}^{n-1}r^t=\frac{1-r^n}{1-r}\ (r\ne1).
\]

\subsection*{Connections}
Induction rests on the successor property of $\N$ (\cref{sec:number-systems}) and proves $\forall n\in\N$ statements
(\cref{sec:quantifiers}). It is used to prove the binomial formula (\cref{ch:binom}) and the factorisation theorem for
polynomials (\cref{ch:func}), and Bernoulli's inequality gives $2^n\ge n+1$, which is useful for limits
(\cref{ch:limits}).

\subsection*{Common mistakes}
\begin{itemize}
  \item Omitting or not actually checking the base.
  \item Not using the hypothesis, or using $P_{n+1}$ (what is to be proved) as if it were known.
  \item Not stating $P_n$ clearly, so that the target $P_{n+1}$ is unclear.
  \item Using too few base cases in the strong form.
  \item Treating verification of many cases as a proof.
\end{itemize}

\subsection*{Exam checklist}
\begin{checklist}
  \item State the principle of induction, naming the base, the step and the hypothesis.
  \item Prove summation identities by induction, marking where the hypothesis is used.
  \item Write the step in the form $P_{n-1}\Rightarrow P_n$.
  \item State the alternative formulation and say when it is useful.
  \item Explain why checking finitely many cases is not a proof.
\end{checklist}

% --------------------------------------------------------------------
\section{Practice questions}
% --------------------------------------------------------------------
\pq[basic] For the claim $\sum_{j=1}^{n}a_j=b_n$ for all $n\in\N$, write down (a) the base, (b) the induction hypothesis,
(c) the statement to be proved in the step.

\pq[mathematical] Prove by induction that $\sum_{j=1}^{n}j^3=\Big(\frac{n(n+1)}2\Big)^2$ for every $n\in\N$.

\pq[mathematical] Prove by induction that $\sum_{j=1}^{n}\frac{1}{j(j+1)}=\frac{n}{n+1}$ for every $n\in\N$.

\pq[mathematical] Prove by induction that $n^3-n$ is divisible by $3$ for every $n\in\N$.

\pq[conceptual] A student proves the step $P_n\Rightarrow P_{n+1}$ for the claim $\sum_{j=1}^{n}(2j-1)=n^2+1$ and concludes
that the claim is true. Explain the error.

\pq[application] A loan of $L$ pounds accrues interest at rate $i$ per year, and nothing is repaid. Let $D_n$ be the debt
after $n$ years, so $D_1=L(1+i)$ and $D_{n+1}=(1+i)D_n$. Prove by induction that $D_n=L(1+i)^n$.

\pq[multi-step] A sequence is defined by $a_1=5$, $a_2=13$ and $a_{n+1}=5a_n-6a_{n-1}$ for $n\ge2$. (a) Explain why ordinary
induction with a single base case is not sufficient, and how many base cases are needed. (b) Prove that $a_n=2^n+3^n$
for every $n\in\N$.

\pq[exam-style] (a) Prove by mathematical induction that $\sum_{j=1}^{n}j(j+1)=\frac{n(n+1)(n+2)}{3}$ for every $n\in\N$.
Identify the base, the hypothesis and the step, and state where the hypothesis is used.
(b) Using (a) and \cref{kf:sum-n}, deduce the formula for $\sum_{j=1}^n j^2$.

% --------------------------------------------------------------------
\section{Solutions to practice questions}
% --------------------------------------------------------------------
\pqsol{7.1}
(a) $a_1=b_1$. (b) For some $n\in\N$, $\sum_{j=1}^{n}a_j=b_n$. (c) $\sum_{j=1}^{n+1}a_j=b_{n+1}$, i.e.\ (after peeling off the
last term) $b_n+a_{n+1}=b_{n+1}$.

\pqsol{7.2}
Base: $1=\big(\frac{1\cdot2}{2}\big)^2=1$. Step: assume the formula for $n$. Then
$\sum_{j=1}^{n+1}j^3=\frac{n^2(n+1)^2}{4}+(n+1)^3=\frac{(n+1)^2\,(n^2+4n+4)}{4}=\frac{(n+1)^2(n+2)^2}{4}
=\Big(\frac{(n+1)(n+2)}2\Big)^2$, which is the formula for $n+1$.

\pqsol{7.3}
Base: $\frac1{1\cdot2}=\frac12=\frac{1}{1+1}$. Step: assume the formula for $n$. Then
$\sum_{j=1}^{n+1}\frac1{j(j+1)}=\frac n{n+1}+\frac1{(n+1)(n+2)}=\frac{n(n+2)+1}{(n+1)(n+2)}=\frac{(n+1)^2}{(n+1)(n+2)}=\frac{n+1}{n+2}$.

\pqsol{7.4}
Base: $1-1=0=3\cdot0$. Step: assume $n^3-n=3k$ for an integer $k$. Then
$(n+1)^3-(n+1)=n^3+3n^2+3n+1-n-1=(n^3-n)+3(n^2+n)=3(k+n^2+n)$, a multiple of $3$.

\pqsol{7.5}
The base fails: for $n=1$ the left side is $1$ and the right side is $2$. The step is valid ($n^2+1+2n+1=(n+1)^2+1$), but
without a true base it establishes nothing --- the correct identity is $\sum(2j-1)=n^2$. Both the base and the step are
needed.

\pqsol{7.6}
Base: $D_1=L(1+i)=L(1+i)^1$. Step: if $D_n=L(1+i)^n$ then $D_{n+1}=(1+i)D_n=L(1+i)^{n+1}$. By induction $D_n=L(1+i)^n$ for all $n$.

\pqsol{7.7}
(a) $a_{n+1}$ is determined by $a_n$ \emph{and} $a_{n-1}$, so the hypothesis $P_n$ alone does not give enough information;
the strong form supplies both $P_n$ and $P_{n-1}$. Two base cases ($n=1$ and $n=2$) are needed, because the recursion
expresses $a_{n+1}$ only for $n\ge2$; $a_2$ cannot be obtained from it (that would need $a_0$).
(b) Base: $a_1=5=2+3$, $a_2=13=4+9$. Step: for $n\ge2$, assume the formula for all indices up to $n$. Then
$a_{n+1}=5(2^n+3^n)-6(2^{n-1}+3^{n-1})=2^{n-1}(10-6)+3^{n-1}(15-6)=4\cdot2^{n-1}+9\cdot3^{n-1}=2^{n+1}+3^{n+1}$.

\pqsol{7.8}
(a) Let $P_n$ be the identity. Base: $1\cdot2=2=\frac{1\cdot2\cdot3}{3}$. Hypothesis: assume $P_n$. Step:
$\sum_{j=1}^{n+1}j(j+1)=\frac{n(n+1)(n+2)}3+(n+1)(n+2)$ (hypothesis used here) $=\frac{(n+1)(n+2)(n+3)}{3}$, which is $P_{n+1}$.
(b) $j^2=j(j+1)-j$, so by P1, P2: $\sum j^2=\frac{n(n+1)(n+2)}3-\frac{n(n+1)}2=\frac{n(n+1)\big(2(n+2)-3\big)}{6}=\frac{n(n+1)(2n+1)}6$.
